\begin{textsample}{POS Dim 1 – 2006 – Score 22.00 – 2006/t000498.txt}  \label{ex:f1_pos_011}
I ' m \textbf{the} luckiest guy in \textbf{the} world.
I got to see \textbf{the} last case of killer smallpox in \textbf{the} world.
I was in India this past year, and I may have seen \textbf{the} last cases of polio in \textbf{the} world.
There ' s nothing that makes you feel more — \textbf{the} blessing and \textbf{the} honor of working in a program like that — than to know that something that horrible no longer exists.
So I ' m going to tell you — so I ' m going to show you some dirty pictures.
They are difficult to watch, but you should look at them with optimism, because \textbf{the} horror of these pictures will be matched by \textbf{the} uplifting quality of knowing that they no longer exist.
But first, I ' m going to tell you a little bit about my own journey.
My background is not exactly \textbf{the} conventional medical education that you might expect.
When I was an intern in San Francisco, I heard about a group of Native Americans who had taken over Alcatraz Island, and a Native American who wanted to give birth on that \textbf{island}, and no other doctor wanted to go and help her give birth.
I went out to Alcatraz, and I lived on \textbf{the} \textbf{island} for several weeks.
She gave birth ; I caught \textbf{the} baby ; I got off \textbf{the} \textbf{island} ; I landed in San Francisco ; and all \textbf{the} press wanted to talk to me, because my three weeks on \textbf{the} \textbf{island} made me an expert in Indian affairs.
I wound up on every television show.
Someone saw me on television ; they called me up ; and they asked me if I ' d like to be in a movie and to play a young doctor for a bunch of \textbf{rock} and roll stars who were traveling in a bus ride from San Francisco to England.
And I said, yes, I would do that, so I became \textbf{the} doctor in an absolutely awful movie called"Medicine Ball Caravan."Now, you know from \textbf{the} ' 60s, you ' re either on \textbf{the} bus or you ' re off \textbf{the} bus ; I was on \textbf{the} bus.
My wife of 37 years and I joined \textbf{the} bus.
Our bus ride took us from San Francisco to London, then we switched buses at \textbf{the} big pond.
We then got on two more buses and we drove through Turkey and Iran, Afghanistan, over \textbf{the} Khyber Pass into Pakistan, like every other young doctor.
This is us at \textbf{the} Khyber Pass, and that ' s our bus.
We had some difficulty getting over \textbf{the} Khyber Pass.
But we wound up in India.
And then, like everyone else in our generation, we went to live in a Himalayan monastery.
This is just like a residency program, for those of you that are in medical school.
And we studied with a wise man, a guru named Karoli Baba, who then told me to get rid of \textbf{the} dress, put on a three-piece suit, go join \textbf{the} United Nations as a diplomat and work for \textbf{the} World Health Organization.
And he made an outrageous prediction that smallpox would be eradicated, and that this was God ' s gift to humanity, because of \textbf{the} hard work of dedicated scientists.
And that prediction came true.
This little girl is Rahima Banu, and she was \textbf{the} last case of killer smallpox in \textbf{the} world.
And this document is \textbf{the} certificate that \textbf{the} global commission signed, certifying \textbf{the} world to have eradicated \textbf{the} first disease in history. \textbf{The} key to eradicating smallpox was early detection, early response.
I ' m going to ask you to repeat that : early detection, early response.
Can you say that?
Audience : Early detection, early response.
Larry Brilliant : Smallpox was \textbf{the} worst disease in history.
It killed more people than all \textbf{the} wars in history.
In \textbf{the} last century, it killed 500 million people.
You ' re reading about Larry Page already.
Somebody reads very fast.
In \textbf{the} year that Larry Page and Sergey Brin — with whom I have a certain affection and a new affiliation — in \textbf{the} year in which they were born, two million people died of smallpox.
We declared smallpox eradicated in 1980.
This is \textbf{the} most important \textbf{slide} that I ' ve ever seen in public health, [ Sovereigns killed by smallpox ] because it shows you to be \textbf{the} richest and \textbf{the} strongest, and to be kings and \textbf{queens} of \textbf{the} world, did not protect you from dying of smallpox.
Never can you doubt that we are all in this together.
But to see smallpox from \textbf{the} perspective of a sovereign is \textbf{the} wrong perspective.
You should see it from \textbf{the} perspective of a mother, watching her child develop this disease and standing by helplessly.
Day one, day two, day three, day four, day five, day six.
You ' re a mother and you ' re watching your child, and on day six, you see pustules that become hard.
Day seven, they show \textbf{the} classic scars of smallpox umbilication.
Day eight.
And Al Gore said earlier that \textbf{the} most photographed image in \textbf{the} world, \textbf{the} most printed image in \textbf{the} world, was that of \textbf{the} Earth.
But this was in 1974, and as of that moment, this photograph was \textbf{the} photograph that was \textbf{the} most widely printed, because we printed two billion \textbf{copies} of this photograph, and we took them hand to hand, door to door, to show people and ask them if there was smallpox in their house, because that was our surveillance system.
We didn ' t have Google, we didn ' t have \textbf{web} crawlers, we didn ' t have \textbf{computers}.
By day nine — you look at this picture and you ' re horrified ; I look at this picture and I say,"Thank God,"because it ' s clear that this is only an ordinary case of smallpox, and I know this child will live.
And by day 13, \textbf{the} lesions are scabbing, his eyelids are swollen, but you know this child has no other secondary infection.
And by day 20, while he will be scarred for life, he will live.
There are other kinds of smallpox that are not like that.
This is confluent smallpox, in which there isn ' t a single place on \textbf{the} body where you could put a finger and not be covered by lesions.
Flat smallpox, which killed 100 percent of people who got it.
And hemorrhagic smallpox, \textbf{the} most cruel of all, which had a predilection for pregnant women.
I ' ve probably had 50 women die.
They all had hemorrhagic smallpox.
I ' ve never seen anybody die from it who wasn ' t a pregnant woman.
In 1967, \textbf{the} WHO embarked on what was an outrageous program to eradicate a disease.
In that year, there were 34 countries affected with smallpox.
By 1970, we were down to 18 countries. 1974, we were down to five countries.
But in that year, smallpox exploded throughout India.
And India was \textbf{the} place where smallpox made its last stand.
In 1974, India had a population of 600 million.
There are 21 linguistic states in India, which is like saying 21 different countries.
There are 20 million people on \textbf{the} road at any time, in buses and trains, walking ; 500, 000 villages, 120 million households, and none of them wanted to report if they had a case of smallpox in their house, because they thought that smallpox was \textbf{the} visitation of a deity, Shitala Mata, \textbf{the} cooling mother, and it was wrong to bring strangers into your house when \textbf{the} deity was in \textbf{the} house.
No incentive to report smallpox.
It wasn ' t just India that had smallpox deities ; smallpox deities were prevalent all over \textbf{the} world.
So, how we eradicated smallpox was — mass vaccination wouldn ' t work.
You could vaccinate everybody in India, but one year later there ' d be 21 million new babies, which was then \textbf{the} population of Canada.
It wouldn ' t do just to vaccinate everyone.
You had to find every single case of smallpox in \textbf{the} world at \textbf{the} same time, and draw a circle of immunity around it.
And that ' s what we did.
In India alone, my 150, 000 best friends and I went door to door, with that same picture, to every single house in India.
We made over one billion house calls.
And in \textbf{the} process, I learned something very important.
Every time we did a house-to-house \textbf{search}, we had a spike in \textbf{the} number of reports of smallpox.
When we didn ' t \textbf{search}, we had \textbf{the} illusion that there was no disease.
When we did \textbf{search}, we had \textbf{the} illusion that there was more disease.
A surveillance system was necessary, because what we needed was early detection, early response.
So we \textbf{searched} and we \textbf{searched}, and we found every case of smallpox in India.
We had a reward.
We raised \textbf{the} reward.
We continued to increase \textbf{the} reward.
We had a scorecard that we wrote on every house.
And as we did that, \textbf{the} number of reported cases in \textbf{the} world dropped to zero.
And in 1980, we declared \textbf{the} globe free of smallpox.
It was \textbf{the} largest campaign in United Nations history, until \textbf{the} Iraq war. 150, 000 people from all over \textbf{the} world — doctors of every race, religion, culture and nation, who fought side by side, brothers and sisters, with each other, not against each other, in a common cause to make \textbf{the} world better.
But smallpox was \textbf{the} fourth disease that was intended for eradication.
We failed three other times.
We failed against \textbf{malaria}, \textbf{yellow} fever and yaws.
But soon we may see polio eradicated.
But \textbf{the} key to eradicating polio is early detection, early response.
This may be \textbf{the} year we eradicate polio.
That will make it \textbf{the} second disease in history.
And David Heymann, who ' s watching this on \textbf{the} webcast — David, keep on going.
We ' re close!
We ' re down to four countries.
I feel like Hank Aaron.
Barry Bonds can replace me any time.
Let ' s get another disease off \textbf{the} list of terrible things to worry about.
I was just in India working on \textbf{the} polio program. \textbf{The} polio surveillance program is four million people going door to door.
That is \textbf{the} surveillance system.
But we need to have early detection, early response.
Blindness, \textbf{the} same thing. \textbf{The} key to discovering blindness is doing epidemiological surveys and finding out \textbf{the} causes of blindness, so you can mount \textbf{the} correct response. \textbf{The} Seva Foundation was started by a group of alumni of \textbf{the} Smallpox Eradication Programme, who, having climbed \textbf{the} highest mountain, tasted \textbf{the} elixir of \textbf{the} success of eradicating a disease, wanted to do it again.
And over \textbf{the} last 27 years, Seva ' s programs in 15 countries have given back sight to more than two million blind people.
Seva got started because we wanted to apply these lessons of surveillance and epidemiology to something which nobody else was looking at as a public health issue : blindness, which heretofore had been thought of only as a clinical disease.
In 1980, Steve Jobs gave me that \textbf{computer}, which is Apple number 12, and it ' s still in Kathmandu, and it ' s still working, and we ought to go get it and auction it off and make more money for Seva.
And we conducted \textbf{the} first Nepal survey ever done for health, and \textbf{the} first nationwide blindness survey ever done, and we had astonishing results.
Instead of finding out what we thought was \textbf{the} case — that blindness was caused mostly by glaucoma and trachoma — we were astounded to find out that blindness was caused instead by cataract.
You can ' t cure or prevent what you don ' t know is there.
In your TED packages there ' s a DVD,"Infinite Vision,"about Dr.
V and \textbf{the} Aravind Eye Hospital.
I hope that you will take a look at it.
Aravind, which started as a Seva project, is now \textbf{the} world ' s largest and best eye hospital.
This year, that one hospital will give back sight to more than 300, 000 people in Tamil Nadu, India.
Bird flu.
I stand here as a representative of all terrible things — this might be \textbf{the} worst. \textbf{The} key to preventing or mitigating pandemic bird flu is early detection and rapid response.
We will not have a vaccine or adequate supplies of an antiviral to combat bird flu if it \textbf{occurs} in \textbf{the} next three years.
WHO stages \textbf{the} progress of a pandemic.
We are now at stage three on \textbf{the} pandemic alert stage, with just a little bit of human-to-human transmission, but no human-to-human-to-human sustained transmission. \textbf{The} moment WHO says we ' ve moved to category four — this will not be like Katrina. \textbf{The} world as we know it will stop.
There ' ll be no \textbf{airplanes} flying.
Would you get in an \textbf{airplane} with 250 people you didn ' t know, coughing and sneezing, when you knew that some of them might carry a disease that could kill you, for which you had no antivirals or vaccine?
I did a study of \textbf{the} top epidemiologists in \textbf{the} world in October.
I asked them — these are all fluologists and specialists in influenza — and I asked them \textbf{the} questions you ' d like to ask them : What do you think \textbf{the} likelihood is that there ' ll be a pandemic?
If it happens, how bad do you think it will be?
Fifteen percent said they thought there ' d be a pandemic within three years.
But much worse than that, 90 percent said they thought there ' d be a pandemic within your children or your grandchildren ' s lifetime.
And they thought that if there was a pandemic, a billion people would get sick.
As many as 165 million people would die.
There would be a global recession and depression as our just-in-time inventory system and \textbf{the} tight rubber band of globalization broke, and \textbf{the} cost to our economy of one to three trillion dollars would be far worse for everyone than merely 100 million people dying, because so many more people would lose their job and their healthcare benefits, that \textbf{the} consequences are almost unthinkable.
And it ' s getting worse, because travel is getting so much better.
Let me show you a simulation of what a pandemic looks like.
So we know what we ' re talking about.
Let ' s assume, for example, that \textbf{the} first case \textbf{occurs} in South Asia.
It initially goes quite slowly.
You get two or three discrete locations.
Then there ' ll be secondary outbreaks, and \textbf{the} disease will spread from country to country so fast that you won ' t know what hit you.
Within three weeks it will be everywhere in \textbf{the} world.
Now, if we had an"undo"button, and we could go back and isolate it and grab it when it first started — if we could find it early, and we had early detection and early response, and we could put each one of those viruses in jail — that ' s \textbf{the} only way to deal with something like a pandemic.
And let me show you why that is.
We have a joke.
This is an epidemic curve, and everyone in medicine, I think, ultimately gets to know what it is.
But \textbf{the} joke is, an epidemiologist likes to arrive at an epidemic right here and ride to glory on \textbf{the} downhill curve.
But you don ' t get to do that usually.
You usually arrive right about here.
What we really want is to arrive right here, so we can stop \textbf{the} epidemic.
But you can ' t always do that.
But there ' s an organization that has been able to find a way to learn when \textbf{the} first cases \textbf{occur}, and that is called GPHIN ; it ' s \textbf{the} Global Public Health Information Network.
And that simulation that I showed you that you thought was bird flu — that was SARS.
And SARS is \textbf{the} pandemic that did not \textbf{occur}.
And it didn ' t \textbf{occur} because GPHIN found \textbf{the} pandemic-to-be of SARS three months before WHO actually announced it, and because of that, we were able to stop \textbf{the} SARS pandemic.
And I think we owe a great debt of gratitude to GPHIN and to Ron St.
John, who I hope is in \textbf{the} audience some place — over there — who ' s \textbf{the} founder of GPHIN.
Hello, Ron!
And TED has flown Ron here from Ottawa, where GPHIN is located, because not only did GPHIN find SARS early, but you may have seen last week that Iran announced that they had bird flu in Iran, but GPHIN found \textbf{the} bird flu in Iran not February 14 — but last September.
We need an early-warning system to protect us against \textbf{the} things that are humanity ' s worst nightmare.
And so my TED wish is based on \textbf{the} common denominator of these experiences.
Smallpox — early detection, early response.
Blindness, polio — early detection, early response.
Pandemic bird flu — early detection, early response.
It is a litany.
It is so obvious that our only way of dealing with these new diseases is to find them early and to kill them before they spread.
So, my TED wish is for you to help build a global system — an early-warning system — to protect us against humanity ' s worst nightmares.
And what I thought I would call it is"Early Detection,"But it should really be called..."Total Early Detection."[ TED ] What?
What?
But in all seriousness, because this idea is birthed in TED, I would like it to be a legacy of TED, and I ' d like to call it \textbf{the}"International System for Total Early Disease Detection."[ INSTEDD ] And INSTEDD then becomes our mantra.
So instead of a hidden pandemic of bird flu, we find it and immediately contain it.
Instead of a novel virus caused by bio-terror or bio-error, or shift or drift, we find it and we contain it.
Instead of industrial accidents like \textbf{oil} \textbf{spills} or \textbf{the} catastrophe in Bhopal, we find them, and we respond to them.
Instead of famine, hidden until it is too late, we \textbf{detect} it, and we respond.
And instead of a system which is owned by a government, and hidden in \textbf{the} bowels of government, let ' s build an early detection system that ' s freely available to anyone in \textbf{the} world in their own language.
Let ' s make it transparent, non-governmental, not owned by any single country or company, housed in a neutral country, with redundant backup in a different time zone and a different \textbf{continent}.
And let ' s build it on GPHIN.
Let ' s start with GPHIN.
Let ' s increase \textbf{the} websites that they crawl from 20, 000 to 20 million.
Let ' s increase \textbf{the} languages they crawl from seven to 70, or more.
Let ' s build in outbound confirmation messages, using text messages or SMS or instant messaging to find out from people who are within 100 meters of \textbf{the} rumor that you hear, if it is, in fact, valid.
And let ' s add satellite confirmation.
And we ' ll add Gapminder ' s amazing \textbf{graphics} to \textbf{the} front end.
And we ' ll grow it as a moral force in \textbf{the} world, finding out those terrible things before anybody else knows about them, and sending our response to them, so that next year, instead of us meeting here, lamenting how many terrible things there are in \textbf{the} world, we will have pulled together, used \textbf{the} unique skills and \textbf{the} magic of this community, and be proud that we have done everything we can to stop pandemics, other catastrophes, and change \textbf{the} world, beginning right now.
Chris Anderson : An amazing presentation.
First of all, just so everyone understands : you ' re saying that by creating \textbf{web} crawlers, looking on \textbf{the} Internet for patterns, they can \textbf{detect} something suspicious before WHO, before anyone else can see it?
Give an example of how that could possibly be true.
Larry Brilliant : You ' re not mad about \textbf{the} copyright violation?
CA : No.
I love it.
LB : Well, as Ron St.
John — I hope you ' ll go and meet him in \textbf{the} dinner afterwards and talk to him.
When he started GPHIN — In 1997, there was an outbreak of bird flu — H5N1.
It was in Hong Kong.
And a remarkable doctor in Hong Kong responded immediately, by slaughtering 1. 5 million chickens and birds, and they stopped that outbreak in its tracks.
Immediate detection, immediate response.
Then a number of years went by, and there were a lot of rumors about bird flu.
Ron and his team in Ottawa began to crawl \textbf{the} \textbf{web} — only crawling 20, 000 different websites, mostly periodicals — and they read about and heard about a concern, of a lot of children who had high fever and symptoms of bird flu.
They reported this to WHO.
WHO took a little while taking action, because WHO will only receive a report from a government, because it ' s \textbf{the} United Nations.
But they were able to point to WHO and let them know that there was this surprising and unexplained cluster of illnesses that looked like bird flu.
That turned out to be SARS.
That ' s how \textbf{the} world found out about SARS.
And because of that, we were able to stop SARS.
Now, what ' s really important is that, before there was GPHIN, 100 percent of all \textbf{the} world ' s reports of bad things — whether you ' re talking about famine or you ' re talking about bird flu or you ' re talking about Ebola — 100 percent of all those reports came from nations. \textbf{The} moment these guys in Ottawa — on a budget of 800, 000 dollars a year — got cracking, 75 percent of all \textbf{the} reports in \textbf{the} world came from GPHIN, 25 percent of all \textbf{the} reports in \textbf{the} world came from all \textbf{the} other 180 nations.
Now, here ' s what ' s really interesting : after they ' d been working for a couple years, what do you think happened to those nations?
They felt pretty stupid.
So they started sending in their reports early.
And now, their reporting percentage is down to 50 percent, because other nations have started to report.
So, can you find diseases early by crawling \textbf{the} \textbf{web}?
Of course you can.
Can you find it even earlier than GPHIN does now?
Of course you can.
You saw that they found SARS using their Chinese \textbf{web} crawler a full six weeks before they found it using their English \textbf{web} crawler.
Well, they ' re only crawling in seven languages.
These bad viruses really don ' t have any intention of showing up first in English or Spanish or French.
So yes, I want to take GPHIN, I want to build on it.
I want to add all \textbf{the} languages of \textbf{the} world that we possibly can.
I want to make this open to everybody, so that \textbf{the} health officer in Nairobi or in Patna, Bihar will have as much access to it as \textbf{the} folks in Ottawa or in CDC.
And I want to make it part of our culture that there is a community of people who are watching out for \textbf{the} worst nightmares of humanity, and that it ' s accessible to everyone.

% matched lemmas: airplane, computer, continent, copy, detect, graphics, island, malaria, occur, oil, queen, rock, search, slide, spill, the, web, yellow
\end{textsample}
